\chapter{Fichiers, scripts et reproductibilité}\label{ann:fichiers}

\begin{tabularx}{\linewidth}{@{}lL@{}}\toprule
Fichier & Contenu\\ \midrule
\code{src/m8\_modele.py} & Modèle macroéconomique v0.8 (fermeture sans termes en différence, θ et ρ estimés, ancrage salarial de l'incidence)\\
\code{src/m8\_reference.py} & Référence calée (recalibrée ; identique sous les deux fermetures à $10^{-13}$ près)\\
\code{src/m15\_offre\_travail.py} & Offre de travail, 18 cas (9 configurations familiales × locataire ou propriétaire)\\
\code{src/m16\_distribution\_dynamique.py} & Distribution dynamique : huit canaux par ménage, cinq horizons, contrôle de cohérence\\
\code{src/m10\_couplage.py} & Microsimulation ; variables par ménage transmises à M16\\
\code{simulateur/executer.py} & Incidence des prélèvements sur les entreprises, surcharges de paramètres, appel de M16\\
\code{phase31/irf\_fermeture.py} & Réponses impulsionnelles de la fermeture (grille ρ, θ, part du capital, amortissement)\\
\code{phase31/estimation\_theta.py} & Durée empirique des écarts de production (AMECO, FMI), sélection de ρ et θ\\
\code{phase31/backtest.py}, \code{analyse\_backtest.py} & Backtest 2010-2025 en temps réel (neuf variantes, cinq origines)\\
\code{phase31/experiences\_v08.py} & 68 expériences relancées avec la v0.8\\
\code{phase31/strategies\_v08.py} & S1-S5 et programmes v0.6, trois configurations\\
\code{phase31/sensibilite.py} & Sensibilité un à un (douze paramètres)\\
\code{phase31/decomposition\_v07\_v08.py} & Décomposition des écarts v0.7/v0.8 (fermeture, incidence)\\
\code{phase31/livrables31.py} & Classeurs Excel de distribution et de comparaison\\
\code{phase31/rapport/generer31.py}, \code{figures\_a.py} & Tables et figures du rapport\\
\code{tests/test\_phase31.py} & Tests automatisés (identités, cohérence, stabilité, reproductibilité, monotonie)\\
\code{doc/18} à \code{doc/20} & Audit v0.7, modèle v0.8, registre des paramètres v0.8\\
\code{data/raw/dbnomics\_*.json} & Séries AMECO et FMI téléchargées (écart de production, potentiel)\\
\code{phase31/resultats/} & Tous les résultats (JSON)\\ \bottomrule
\end{tabularx}

\paragraph{Ordre d'exécution.} \code{phase3/caler\_reference.py} puis \code{phase3/maj\_cache\_reference.py} (environ
45 min) ; \code{estimation\_theta.py} ; \code{backtest.py} et \code{analyse\_backtest.py} ; \code{experiences\_v08.py
lots …} (environ 70 min) ; \code{strategies\_v08.py central faible fort} ; \code{sensibilite.py} ;
\code{decomposition\_v07\_v08.py} ; \code{phase3/optimisation.py} avec \code{RES\_OPTIM} et \code{DISTRIB\_DYN} ;
\code{livrables31.py} ; \code{rapport/generer31.py} ; \code{pytest tests/test\_phase31.py} ; compilation XeLaTeX.

\paragraph{Traçabilité.} Les résultats v0.6 (\code{rapport/resultats/}) et v0.7 (\code{phase3/resultats/}) ne sont
pas modifiés. Les réponses impulsionnelles calculées avant la correction du conflit de noms (\code{irf\_v1}) et avec
μ = 0,3 (\code{irf\_mu03}) sont conservées pour mémoire.

\chapter{Tests automatisés}\label{ann:tests}

{\small\begin{longtable}{@{}p{5.2cm}p{8.2cm}cr@{}}\toprule
Test & Propriété vérifiée & Résultat & s \\ \midrule\endhead
\code{test\_\allowbreak{}dette\_\allowbreak{}et\_\allowbreak{}solde} & Identité : dette(t) = dette(t−1) − solde(t) (hors ajustements stock-flux de calage, nuls après 2026). & ok & 11 \\
\code{test\_\allowbreak{}ratios} & Les ratios au PIB sont cohérents avec les niveaux ; solde = recettes − dépenses. & ok & 0 \\
\code{test\_\allowbreak{}reference\_\allowbreak{}identique\_\allowbreak{}sous\_\allowbreak{}les\_\allowbreak{}deux\_\allowbreak{}fermetures} & reference identique sous les deux fermetures & ok & 4 \\
\code{test\_\allowbreak{}reproductibilite} & Deux simulations identiques donnent des résultats identiques (aucun aléa caché). & ok & 4 \\
\code{test\_\allowbreak{}amortissement\_\allowbreak{}du\_\allowbreak{}point\_\allowbreak{}fixe\_\allowbreak{}independant\_\allowbreak{}de\_\allowbreak{}la\_\allowbreak{}fermeture} & Le défaut de la v0.7 (le terme proportionnel de la fermeture écrasait l'amortissement du point fixe) ne doit pas revenir : changer le terme proportionnel à partir de 2028 ne modifie pas les années antérieures. & ok & 5 \\
\code{test\_\allowbreak{}choc\_\allowbreak{}ponctuel\_\allowbreak{}sans\_\allowbreak{}oscillation} & Après un choc d'offre ponctuel (productivité −1 \% en 2028), l'écart de production décroît sans changer de signe de plus de 10 \% de l'écart initial (pas d'oscillation de période deux ans). & ok & 5 \\
\code{test\_\allowbreak{}chocs\_\allowbreak{}d\_\allowbreak{}offre} & Un choc de productivité ou de démographie défavorable réduit le PIB et le PIB potentiel ; l'écart de production reste faible (pas de comportement paradoxal). & ok & 9 \\
\code{test\_\allowbreak{}taux} & Une hausse des taux réduit l'investissement et augmente la dette. & ok & 5 \\
\code{test\_\allowbreak{}incidence\_\allowbreak{}is\_\allowbreak{}sur\_\allowbreak{}les\_\allowbreak{}salaires} & Une hausse d'IS réduit le coût salarial de long terme quand l'incidence salariale est positive, pas sinon. & ok & 86 \\
\code{test\_\allowbreak{}distribution\_\allowbreak{}dynamique\_\allowbreak{}coherente} & Le contrôle de cohérence de M16 : la somme des canaux micro reproduit la variation macroéconomique du revenu disponible à 1,5 point près ; le canal direct est nul avant la mise en œuvre (2027). & ok & 42 \\
\code{test\_\allowbreak{}distribution\_\allowbreak{}dynamique\_\allowbreak{}sens} & Une hausse du RSA profite d'abord au premier décile ; une hausse de TVA réduit le revenu réel moyen. & ok & 40 \\
\code{test\_\allowbreak{}monotonie\_\allowbreak{}offre\_\allowbreak{}de\_\allowbreak{}travail} & Les effets de la prime d'activité sur la population active sont ordonnés selon les élasticités basse, centrale, haute. & ok & 0 \\
\code{test\_\allowbreak{}monotonie\_\allowbreak{}anticipation} & Plus l'anticipation est forte, plus la production rejoint vite le potentiel après un choc d'offre. & ok & 13 \\
\code{test\_\allowbreak{}distribution\_\allowbreak{}reforme\_\allowbreak{}age\_\allowbreak{}et\_\allowbreak{}assurance\_\allowbreak{}chomage} & Correction phase 3.1 : le relèvement de l'âge légal ne fait pas du premier décile un gagnant (les emplois maintenus vont aux ménages proches de l'âge légal, la pension non versée est une perte) ; la baisse des allocations chômage fait perdre d'abord les premiers déciles ; la cohérence micro-macro est assurée. & ok & 105 \\
\bottomrule\end{longtable}}

\noindent Bilan : 14 tests réussis, 0 en échec.

